在模态分配问题中会出现 $\binom Sk^2$ 型权重。要讨论它的峰宽，先把归一化与随机变量的取样含义算清楚。

\begin{example}[平方二项式权重的超几何分布]
设 $S\ge1$ 为整数，总体有 $2S$ 个对象，其中两类各有 $S$ 个；不放回等概率抽取 $S$ 个，令 $K$ 为第一类的个数，则
\begin{equation}\label{eq:math-ghs-hypergeometric-probability}
 P(K=k)=\frac{\binom Sk\binom S{S-k}}{\binom{2S}S}
 =\frac{\binom Sk^2}{\binom{2S}S},\qquad 0\le k\le S.
\end{equation}
分子分别选择两类对象，分母计数所有子集。Vandermonde 恒等式 $\sum_k\binom Sk\binom S{S-k}=\binom{2S}S$ 直接验证归一化。

令 $X_i$ 指示第 $i$ 个第一类对象是否入选。每个 $P(X_i=1)=1/2$，而不同对象满足
\[
 P(X_i=X_j=1)=\frac{S(S-1)}{2S(2S-1)}
 =\frac{S-1}{2(2S-1)}.
\]
因此 $\operatorname{Cov}(X_i,X_j)=-1/[4(2S-1)]$。将 $K=\sum_{i=1}^SX_i$ 代入，得到
\begin{equation}\label{eq:math-ghs-hypergeometric-moments}
 \mathbb E K=\frac S2,\qquad
 \operatorname{Var}K=\frac S4-\frac{S(S-1)}{4(2S-1)}
 =\frac{S^2}{4(2S-1)}.
\end{equation}
不放回抽样产生负相关，故方差小于独立二项分配的 $S/4$。当 $S\to\infty$ 时，标准差为 $\sqrt{S/8}[1+O(S^{-1})]$；权重围绕中点集中，但其绝对宽度仍随 $\sqrt S$ 增长。

\end{example}

《自由电子量子光学》的高自旋双粒子态给出同一权重的模式占据来源。这里建立的是抽样分布的数学性质；态归一化、单体约化谱和熵必须结合该书的 Fock 空间定义，不能仅从概率形式推断。
